\chapter{Research Chronology}
\label{app:chronology}
\section{Original Objective and Dataset Selection}
The original research objective concerned process-aware multimodal architectures for robotic welding. The exposé defined the research questions before the subsequent industrial and robotic studies. The choice and timing of those studies were exploratory. Their relationship to the original questions is a retrospective interpretation.\ev{Original scope and scientific chronology}{\EvidenceRef{E002}{pp. 2--3}, \EvidenceRef{E004}{sections 10--11}}

\section{From the Welding Motivation to Two Task-Specific Studies}
The PHM study began with image exploration and developed a continuous gear-condition proxy. The September v3 target used one deterministic view per tooth, teeth 5--28 and a top-three run aggregate. The earlier maximum-over-views targets and fixed image split addressed different prediction settings. The available records do not establish the initial reason for selecting PHM. Concerns about the validity of the proxy emerged through its construction and evaluation.\ev{PHM target chronology}{\EvidenceRef{E004}{sections 12--14}, \EvidenceRef{E023}{}}

Re-examination of the Intel welding data identified shared sessions across partitions and prior access to the proposed test set. No independent multimodal evaluation of weld quality followed. PHM and REASSEMBLE provided settings for the architecture and evaluation questions. In REASSEMBLE, native action outcomes and corresponding camera and telemetry observations supported comparison of the fusion methods. The experiment excluded audio because annotation times could not be reliably mapped to waveform times and spoken outcomes could reveal the labels.\ev{Welding and REASSEMBLE transition}{\EvidenceRef{E003}{sections 2--3}, \EvidenceRef{E004}{sections 14--15}}

\section{Follow-Ups, Human Review and External Cases}
September follow-ups examined longer sensor training, a same-sensor hybrid, anchored adaptive coefficients, matched augmentation and boundary-aware visual pooling. These include positive exploratory results, but they did not execute a combined system containing every favorable choice. The discontinued temporal-model work is distinct from completed pooling experiments.

The October experiments examined a spatial hybrid encoder, fine-tuning with unequal training budgets, a matched two-epoch comparison and noise/quality interactions on selected folds. GPU availability limited the 100-epoch comparison to one previously inspected assessment fold. The later F2 fit also used a different set of training predictions from the original experiment. The results therefore describe a follow-up on that fold.\ev{Follow-up chronology and scope}{\EvidenceRef{E004}{sections 18--21}, \EvidenceRef{E003}{sections 3.3 and 4.4}}

A reviewer used the September interface to assess 64 REASSEMBLE clips. The returned ratings identified disagreements and potential boundary or presentation problems. Adjudication was incomplete, and the evaluation used the original labels.

The October MSM feasibility subset contained too few groups with positive cases. A subsequent evaluation included 1,895 intervals and varied both the grouping and the treatment of native label~2. HANDOVER evaluated five baselines on common input windows. Missing participant identities prevented participant-disjoint fitting of learned fusion. ViFailback progressed from cohort and telemetry checks to an evaluation of the original frozen backbone and a separate matched study of donor weights.\ev{Later external-study reconciliation}{\EvidenceRef{E003}{sections 3.4--3.6 and 4.5--4.7}, \EvidenceRef{E108}{}}

